\section{Analysis of Model~$B_2$ ($\gamma_1>0$, $\gamma_2=0$, $\theta_1=\theta_2=0$)}\label{sec:model_b2}

Model~$B_2$ restricts jockeying to the direction $1\to2$ only ($\gamma_1>0$, $\gamma_2=0$) and retains no abandonment ($\theta_1=\theta_2=0$). Setting $\gamma_2=0$ eliminates the $\partial P/\partial y$ term from the fundamental equation, reducing it to a first-order linear ODE in $x$ for each fixed $y\in[0,1]$. An integrating factor then yields a formal integral expression for $P(x,y)$; the boundary function $P_y(y)$ is pinned by requiring $P(\cdot,y)$ to be holomorphic at the interior singular point $x=y$, where the integrating factor carries a branch singularity with negative exponent. This analyticity condition replaces the boundedness argument used for Model~$C_2$.

\textbf{Stability.} As in Models~$A$ and~$B$, $\rho = (\lambda_1+\lambda_2)/\mu < 1$. Jockeying conserves the total customer count and does not affect the stability condition.

\begin{figure}[H]
    \centering
    \resizebox{0.62\textwidth}{!}{%
        % Model B₂: γ₂=0, θ₁=θ₂=0, γ₁>0
% One-way jockeying: Class-1 can jockey into Class-2 queue only (dashed, downward).
% Class-2 customers cannot jockey back.
\begin{tikzpicture}[>=Latex, thick]

    % ── Queue 1 (top, Class-1, priority) ──────────────────────────────────────
    \draw (0, 1.9) -- (2, 1.9) -- (2, 1.1) -- (0, 1.1);
    \foreach \x in {0.2, 0.4, ..., 1.8} { \draw[thin] (\x, 1.8) -- (\x, 1.2); }
    \draw[->] (-1.5, 1.5) -- (-0.1, 1.5) node[midway, above] {$\lambda_1$};

    % ── Queue 2 (bottom, Class-2) ──────────────────────────────────────────────
    \draw (0, -1.1) -- (2, -1.1) -- (2, -1.9) -- (0, -1.9);
    \foreach \x in {0.2, 0.4, ..., 1.8} { \draw[thin] (\x, -1.2) -- (\x, -1.8); }
    \draw[->] (-1.5, -1.5) -- (-0.1, -1.5) node[midway, above] {$\lambda_2$};

    % ── One-way jockeying (1→2 only) ─────────────────────────────────────────
    \draw[->, dashed] (1.0, 1.0) -- (1.0, -1.0) node[midway, left] {$\gamma_1$};
    % No upward γ₂ arrow

    % ── Server ────────────────────────────────────────────────────────────────
    \node[draw, circle, minimum size=1.2cm] (server) at (5, 0) {$\mu$};
    \draw[->] (2.2,  1.5) -- (server);
    \draw[->] (2.2, -1.5) -- (server);
    \draw[->] (server) -- (7, 0);

\end{tikzpicture}
%
    }
    \caption{Queue layout for Model~$B_2$. Only the $1\to2$ jockeying direction
    is active (dashed, downward): Class-1 customers may defect to Queue~2,
    but no customer can move in the reverse direction ($\gamma_2=0$).}
    \label{fig:model-b2-queue}
\end{figure}

\subsection{Reduced fundamental equation}\label{ssec:B2:reduction}

Setting $\gamma_2=0$ and $\theta_1=\theta_2=0$ in the general fundamental
equation~\eqref{eq:gen:fundamental_equation} (Lemma~\ref{lem:gen:fundamental_equation})
annihilates the $\partial P/\partial y$ term, whose coefficient
$xy\bigl[\gamma_2(y-x)+\theta_2(y-1)\bigr]$ vanishes identically, and leaves only the
$\gamma_1$ contribution to the $\partial P/\partial x$ term:
\begin{align}
    &\bigl[(\lambda_1+\lambda_2+\mu)xy - \mu y - \lambda_1 x^2 y - \lambda_2 xy^2\bigr]P
    + \gamma_1 xy(x-y)\,\frac{\partial P}{\partial x} \nonumber \\
    &\qquad = \mu(x-y)\,P_y(y) - \mu x(1-y)\,\pi(0,0).
    \label{eq:B2:fundamental}
\end{align}
For each fixed $y\in(0,1]$ this is a first-order linear ODE in $x$; the entire analysis
of this section is built on it.

\begin{theorem}\label{thm:B2}
    Consider the two-class non-preemptive priority queueing system of Model~$B_2$ with
    per-customer jockeying rate $\gamma_1>0$ from Queue~1 to Queue~2 ($\gamma_2=0$,
    $\theta_1=\theta_2=0$). Under the stability condition $\rho=\lambda/\mu<1$,
    $\lambda=\lambda_1+\lambda_2$, the boundary generating function is
    \begin{equation}
        P_y(y) = \frac{\mu\,\pi(0,0)}{\mu-\lambda y}
        \cdot\frac{{}_1F_1\!\bigl(\alpha(y)+1;\,b^\star(y);\,-\lambda_1 y/\gamma_1\bigr)}
                  {{}_1F_1\!\bigl(\alpha(y);\,b^\star(y);\,-\lambda_1 y/\gamma_1\bigr)},
        \label{eq:B2:Py_thm}
    \end{equation}
    and the bivariate PGF, for $0\le x\le y$, is
    \begin{align}
        P(x,y) &= \frac{\mu\,e^{\lambda_1 x/\gamma_1}\,x^{-\alpha}(y-x)^{-\beta}}{\gamma_1 y}
        \Bigl[P_y(y)\,\mathcal{I}_1(x,y)
              + (1-y)\pi(0,0)\,\mathcal{I}_2(x,y)\Bigr],
        \label{eq:B2:Pxy}
    \end{align}
    where ${}_1F_1$ is Kummer's confluent hypergeometric function, the exponents are
    \begin{equation}
        \alpha(y) = \frac{\mu}{\gamma_1 y}, \qquad
        \beta(y)  = \frac{(1-y)(\lambda y-\mu)}{\gamma_1 y}, \qquad
        b^\star(y) = \alpha+\beta+1 = \frac{\mu+\lambda(1-y)+\gamma_1}{\gamma_1},
        \label{eq:B2:exponents_thm}
    \end{equation}
    the auxiliary integrals are
    \begin{align}
        \mathcal{I}_1(x,y) &= \int_0^x e^{-\lambda_1 t/\gamma_1}\,t^{\,\alpha-1}(y-t)^{\beta}\,dt,
        \label{eq:B2:I1}\\
        \mathcal{I}_2(x,y) &= \int_0^x e^{-\lambda_1 t/\gamma_1}\,t^{\,\alpha}(y-t)^{\beta-1}\,dt,
        \label{eq:B2:I2}
    \end{align}
    and the empty-state probability $\pi(0,0)=\rho(1-\rho)$ is given by
    Corollary~\ref{cor:B2:pi00} below.
    On $y<x\le1$ the same expression~\eqref{eq:B2:Pxy} holds with $\mathcal{I}_1,\mathcal{I}_2$
    replaced by their finite-part continuations.
\end{theorem}

\begin{corollary}\label{cor:B2:pi00}
    The empty-system probability $\pi_0$ (server idle, no one waiting) and the
    both-queues-empty probability $\pi(0,0)$ (server busy, $N_1=N_2=0$) of Model~$B_2$ are
    \begin{equation}
        \pi_0 = 1-\rho, \qquad \pi(0,0) = \rho(1-\rho),
        \label{eq:B2:pi00}
    \end{equation}
    identical to the Model~$A$ baseline.
\end{corollary}

\begin{proof}
    Because one-way jockeying conserves the total customer count and there is no
    abandonment, the total-count process $L$ (equal to the number in queue $N$ plus the
    single customer in service whenever the server is busy) is itself a birth--death
    chain: it increases at rate $\lambda=\lambda_1+\lambda_2$ on any arrival and decreases
    at rate $\mu$ on any service completion, irrespective of class. Thus $L$ is an $M/M/1$
    queue, and by the memorylessness of the exponential clocks its stationary law is
    geometric, $\mathbb{P}(L=k)=(1-\rho)\rho^k$; in particular
    $\pi_0=\mathbb{P}(L=0)=1-\rho$.

    For $\pi(0,0)$, set $x=y$ in the reduced fundamental
    equation~\eqref{eq:B2:fundamental}. The convection term $\gamma_1 xy(x-y)\partial_x P$
    and the boundary term $\mu(x-y)P_y(y)$ both carry the factor $(x-y)$ and vanish on the
    diagonal; dividing the surviving terms by the common factor $x$ leaves
    \begin{equation}
        \bigl[(\lambda_1+\lambda_2+\mu)x - \mu - (\lambda_1+\lambda_2)x^2\bigr]P(x,x)
        = -\mu(1-x)\,\pi(0,0).
        \label{eq:B2:diagonal_ode}
    \end{equation}
    Factoring the bracket and cancelling $(x-1)$ against $-(1-x)=(x-1)$ on the right:
    \begin{align}
        \bigl[{-\lambda x^2} + (\lambda+\mu)x - \mu\bigr]P(x,x)
        &= (\mu-\lambda x)(x-1)\,P(x,x) \nonumber\\
        &= \mu(x-1)\,\pi(0,0),
        \label{eq:B2:diagonal_factored}
    \end{align}
    so that
    \begin{equation}
        P(x,x) = \frac{\mu\,\pi(0,0)}{\mu-\lambda x} = \frac{\pi(0,0)}{1-\rho x}.
        \label{eq:B2:Pxx}
    \end{equation}
    Evaluating at $x=1$ with $P(1,1)=1-\pi_0=\rho$ (using $\pi_0=1-\rho$ just established)
    yields $\pi(0,0)=\rho(1-\rho)$.
\end{proof}

\begin{proof}[Proof of Theorem~\ref{thm:B2}]

% ─────────────────────────────────────────────────────────────────
\noindent\textit{Stage~1: Standard form and the integrating factor.}
\medskip

For any fixed $y\neq0$, divide the reduced fundamental
equation~\eqref{eq:B2:fundamental} through by the coefficient $\gamma_1 xy(x-y)$ of
$\partial P/\partial x$. The factor $y$ cancels against the overall factor $y$ in the
bracket multiplying $P$, putting the equation in standard first-order form
$\tfrac{dP}{dx}+p(x;y)\,P=q(x;y)$ with
\begin{align}
    p(x;y) &= \frac{(\lambda_1+\lambda_2+\mu)x - \mu - \lambda_1 x^2 - \lambda_2 xy}
                   {\gamma_1 x(x-y)},
    \label{eq:B2:pxy}\\[4pt]
    q(x;y) &= \frac{\mu\,P_y(y)}{\gamma_1 xy}
             - \frac{\mu(1-y)\,\pi(0,0)}{\gamma_1 y(x-y)}.
    \label{eq:B2:qxy}
\end{align}
Write the numerator of $p$ as a polynomial in $x$,
\[
    N(x) = -\lambda_1 x^2 + \bigl[(\lambda_1+\lambda_2+\mu)-\lambda_2 y\bigr]x - \mu
         = -\lambda_1\bigl(x-x^*(y)\bigr)\bigl(x-x^+(y)\bigr),
\]
which factors through the roots $x^*(y),x^+(y)$ of the Model~$A$ kernel
(Equation~\eqref{eq:A:xstar}), with $\lambda_1 x^* x^+ = \mu$. The partial-fraction
coefficients of $p = N(x)/\bigl[\gamma_1 x(x-y)\bigr]$ follow by residues:
\begin{align}
    A &= -\lambda_1 && \text{(ratio of leading $x^2$ terms)}, \nonumber\\
    B &= \frac{N(0)}{0-y} = \frac{-\mu}{-y} = \frac{\mu}{y}, \nonumber\\
    C &= \frac{N(y)}{y} = \frac{-(\lambda y-\mu)(y-1)}{y}
       = \frac{(1-y)(\lambda y-\mu)}{y},
    \label{eq:B2:pfrac_coeffs}
\end{align}
where $N(y)=-\lambda y^2+(\lambda+\mu)y-\mu=-(\lambda y-\mu)(y-1)$. Hence
\begin{equation}
    p(x;y) = \frac{1}{\gamma_1}\!\left[
        -\lambda_1 + \frac{\mu/y}{x} + \frac{(1-y)(\lambda y-\mu)/y}{x-y}
    \right].
    \label{eq:B2:pfrac}
\end{equation}
Integrating term by term yields the integrating factor
\begin{equation}
    \mu_I(x;y) = e^{-\lambda_1 x/\gamma_1}\cdot x^{\,\alpha(y)}\cdot(x-y)^{\,\beta(y)},
    \label{eq:B2:mu_I}
\end{equation}
with exponents $\alpha(y)=\mu/(\gamma_1 y)>0$ and
$\beta(y)=(1-y)(\lambda y-\mu)/(\gamma_1 y)$ as in~\eqref{eq:B2:exponents_thm}.
Since $\alpha>0$, $\mu_I\to0$ as $x\to0^+$, so the boundary term $P(0,y)\mu_I(0;y)$
vanishes when we integrate $\tfrac{d}{dx}[P\cdot\mu_I]=q\cdot\mu_I$ over $(0,x]$:
\begin{equation}
    P(x,y)\,\mu_I(x;y) = \int_0^x q(t;y)\,\mu_I(t;y)\,dt.
    \label{eq:B2:ODE_integral}
\end{equation}

% ─────────────────────────────────────────────────────────────────
\medskip\noindent\textit{Stage~2: Determination of $P_y(y)$ by analyticity at $x=y$.}
\medskip

Unlike Models~$A$ and~$C_2$, no probabilistic (Pollaczek--Khinchine) route is available
here. One-way jockeying makes the effective service experienced by a Class-2 customer
depend on the Class-1 queue contents at the (random) jockeying instants, so the Class-2
input is no longer a Poisson stream feeding an $M/G/1$ subsystem and no busy-period
reduction applies (cf.\ the discussion of Model~$B$ in \S\ref{sec:model_b}). We therefore
determine $P_y(y)$ \emph{purely analytically}, by imposing holomorphy of $P(\cdot,y)$ at
the interior point $x=y$.

Under stability, $\mu-\lambda y>0$ for all $y\in(0,1)$, so
\begin{equation}
    \beta(y) = \frac{(1-y)(\lambda y-\mu)}{\gamma_1 y} < 0,
    \qquad 0<y<1.
    \label{eq:B2:beta_negative}
\end{equation}
Hence $(x-y)^{\beta}\to\infty$ as $x\to y$: the integrating factor $\mu_I$ diverges at the
interior point $x=y$, and the solution $P=\mu_I^{-1}\!\int q\,\mu_I\,dt$ is a $0\cdot\infty$
form there that must be resolved.  The key observation is that $P(\cdot,y)$, being a
probability generating function, must be holomorphic on the open unit disc; the
non-integer-power branch term $(x-y)^{|\beta|}$ in the general solution must therefore
vanish.

\smallskip\noindent\textit{Branch resolution.}
For $0<t<x<y$ the base $(t-y)<0$, so $(t-y)^\beta = e^{i\pi\beta}(y-t)^\beta$ and
$(t-y)^{\beta-1}=e^{i\pi(\beta-1)}(y-t)^{\beta-1}$.
Substituting these into $P=\mu_I^{-1}\int_0^x q\,\mu_I\,dt$, the common phase $e^{i\pi\beta}$
shared by $\mu_I^{-1}(x)\propto(x-y)^{-\beta}$ and the $\mathcal{I}_1$-integrand cancels,
while the $\mathcal{I}_2$-integrand carries the \emph{relative} phase
$e^{i\pi(\beta-1)}/e^{i\pi\beta}=e^{-i\pi}=-1$, which converts the minus sign
in~\eqref{eq:B2:qxy} into a plus.  The manifestly real form of the
solution~\eqref{eq:B2:ODE_integral} on $0<x<y$ is therefore
\begin{equation}
    P(x,y) = \frac{\mu\,e^{\lambda_1 x/\gamma_1}\,x^{-\alpha}(y-x)^{-\beta}}{\gamma_1 y}
    \Bigl[P_y(y)\,\mathcal{I}_1(x,y) + (1-y)\pi(0,0)\,\mathcal{I}_2(x,y)\Bigr],
    \label{eq:B2:Pxy_real}
\end{equation}
with $\mathcal{I}_1, \mathcal{I}_2$ as defined in~\eqref{eq:B2:I1}--\eqref{eq:B2:I2}.

\smallskip\noindent\textit{Local analysis at $x=y$.}
The ODE~\eqref{eq:B2:pxy}--\eqref{eq:B2:qxy} has a simple pole at $x=y$ with residue $\beta$
in $p$ and a leading singularity $-D/(x-y)$ in $q$, where
$D=\mu(1-y)\pi(0,0)/(\gamma_1 y)$.  The local model equation is
\begin{equation}
    P'(x,y) + \frac{\beta}{x-y}\,P(x,y)
    = \frac{-D}{x-y} + (\text{analytic near }x=y).
    \label{eq:B2:local_ode}
\end{equation}
A constant particular solution $P\equiv P_*$ requires $\beta P_* = -D$, giving
\begin{equation}
    P(y,y) = P_* = -\frac{D}{\beta}
    = \frac{\mu\,\pi(0,0)}{\mu-\lambda y}
    = \frac{\rho(1-\rho)}{1-\rho y},
    \label{eq:B2:regular_value}
\end{equation}
consistent with~\eqref{eq:B2:Pxx} (using $\pi(0,0)=\rho(1-\rho)$). The general
local solution adds a multiple of the homogeneous mode
$(x-y)^{-\beta}=(x-y)^{|\beta|}$. Since $|\beta|\notin\mathbb{Z}$ generically, this
is a branch term: it vanishes at $x=y$ but is not analytic there. Holomorphy therefore
forces its coefficient to zero.

\smallskip\noindent\textit{No-branch condition and $P_y(y)$.}
As $x\to y^-$ the two integrals in~\eqref{eq:B2:Pxy_real} each split into a divergent part
and a finite part.  The divergent part of $\mathcal{I}_2$ (scaling like
$(y-x)^{\beta}$) combines with the prefactor $(y-x)^{-\beta}$ to reproduce exactly the
regular value~\eqref{eq:B2:regular_value}, and the divergent part of $\mathcal{I}_1$
contributes an integer power $(y-x)^{1}$; both are analytic at $x=y$.  The only
non-analytic contribution is the prefactor $(y-x)^{-\beta}=(y-x)^{|\beta|}$ multiplying the
\emph{finite parts} of $\mathcal{I}_1,\mathcal{I}_2$.  Eliminating this branch term requires
\begin{equation}
    P_y(y)\,\mathrm{f.p.}\,\mathcal{I}_1(y,y)
    + (1-y)\,\pi(0,0)\,\mathrm{f.p.}\,\mathcal{I}_2(y,y) = 0,
    \label{eq:B2:nobranch}
\end{equation}
where $\mathrm{f.p.}$ denotes the Hadamard finite part (the analytic continuation in the
exponent parameter).  Substituting $t=ys$ in~\eqref{eq:B2:I1}--\eqref{eq:B2:I2} gives
$\int_0^1 s^{\alpha-1}(1-s)^{\beta}e^{-cs}\,ds$ (resp.\ $s^{\alpha}(1-s)^{\beta-1}$) with
$c=\lambda_1 y/\gamma_1$, to which we apply the Euler integral
representation~\eqref{eq:kummer:euler}, read as its analytic continuation in the second
parameter $b$ beyond the convergence domain $b>0$ (here $b=\beta+1$, resp.\ $b=\beta$, with
$\beta<0$ by~\eqref{eq:B2:beta_negative} --- this is precisely the finite part):
\begin{align}
    \mathrm{f.p.}\,\mathcal{I}_1(y,y)
        &= B(\alpha,\,\beta+1)\,y^{\alpha+\beta}\,
           {}_1F_1\!\bigl(\alpha;\,b^\star(y);\,-\lambda_1 y/\gamma_1\bigr),
    \label{eq:B2:fp1}\\[3pt]
    \mathrm{f.p.}\,\mathcal{I}_2(y,y)
        &= B(\alpha+1,\,\beta)\,y^{\alpha+\beta}\,
           {}_1F_1\!\bigl(\alpha+1;\,b^\star(y);\,-\lambda_1 y/\gamma_1\bigr).
    \label{eq:B2:fp2}
\end{align}
The common factor $y^{\alpha+\beta}$ cancels in~\eqref{eq:B2:nobranch}.  The Beta ratio
reduces, via $\Gamma(\alpha+1)=\alpha\Gamma(\alpha)$ and $\Gamma(\beta+1)=\beta\Gamma(\beta)$,
to
\begin{align}
    \frac{B(\alpha+1,\,\beta)}{B(\alpha,\,\beta+1)}
    &= \frac{\Gamma(\alpha+1)\Gamma(\beta)}{\Gamma(\alpha+\beta+1)}
       \cdot\frac{\Gamma(\alpha+\beta+1)}{\Gamma(\alpha)\Gamma(\beta+1)} \nonumber\\
    &= \frac{\Gamma(\alpha+1)}{\Gamma(\alpha)}\cdot\frac{\Gamma(\beta)}{\Gamma(\beta+1)}
       = \frac{\alpha}{\beta},
    \label{eq:B2:beta_ratio}
\end{align}
and the prefactor combination simplifies, using
$\alpha/\beta=\mu/\bigl[(1-y)(\lambda y-\mu)\bigr]$, as
\begin{align}
    (1-y)\left(-\frac{\alpha}{\beta}\right)
    &= -(1-y)\cdot\frac{\mu}{(1-y)(\lambda y-\mu)} \nonumber\\
    &= -\frac{\mu}{\lambda y-\mu} = \frac{\mu}{\mu-\lambda y}.
    \label{eq:B2:prefactor_combo}
\end{align}
Solving~\eqref{eq:B2:nobranch} for $P_y(y)$:
\begin{align}
    P_y(y) &= \frac{\mu\,\pi(0,0)}{\mu-\lambda y}
    \cdot\frac{{}_1F_1\!\bigl(\alpha(y)+1;\,b^\star(y);\,-\lambda_1 y/\gamma_1\bigr)}
              {{}_1F_1\!\bigl(\alpha(y);\,b^\star(y);\,-\lambda_1 y/\gamma_1\bigr)},
    \label{eq:B2:Py}
\end{align}
confirming~\eqref{eq:B2:Py_thm}.  The prefactor equals the regular value $P(y,y)$
of~\eqref{eq:B2:regular_value}, and the Kummer ratio is the boundary correction.
As $y\to0^+$ the argument $-\lambda_1 y/\gamma_1\to0$, so the ratio tends to $1$ and
$P_y(0)=\pi(0,0)$, consistent with
$P_y(0)=\sum_{n_2}\pi(0,n_2)\cdot0^{n_2}=\pi(0,0)$.

% ─────────────────────────────────────────────────────────────────
\medskip\noindent\textit{Stage~3: Recovery of $P(x,y)$.}
\medskip

Substituting~\eqref{eq:B2:Py} into the real-form expression~\eqref{eq:B2:Pxy_real}
gives, for $0\le x\le y$,
\begin{align}
    P(x,y) &= \frac{\mu\,e^{\lambda_1 x/\gamma_1}\,x^{-\alpha}(y-x)^{-\beta}}{\gamma_1 y}
    \biggl[
        \frac{\mu\,\pi(0,0)}{\mu-\lambda y}
        \cdot\frac{{}_1F_1\!\bigl(\alpha+1;\,b^\star;\,-\tfrac{\lambda_1 y}{\gamma_1}\bigr)}
                  {{}_1F_1\!\bigl(\alpha;\,b^\star;\,-\tfrac{\lambda_1 y}{\gamma_1}\bigr)}
        \cdot\mathcal{I}_1(x,y)
    \nonumber\\[4pt]
    &\hspace{5.8cm}
        + (1-y)\pi(0,0)\cdot\mathcal{I}_2(x,y)
    \biggr],
    \label{eq:B2:Pxy_proof}
\end{align}
which is the statement of the theorem (Equation~\eqref{eq:B2:Pxy}).
The no-branch condition~\eqref{eq:B2:nobranch} guarantees that this extends
analytically across $x=y$, with diagonal value
$P(y,y)=\rho(1-\rho)/(1-\rho y)$ as computed in~\eqref{eq:B2:regular_value}.
On $y<x\le1$ the same formula holds with $\mathcal{I}_1,\mathcal{I}_2$ read as their
finite-part continuations; $\pi(0,0)=\rho(1-\rho)$ from Corollary~\ref{cor:B2:pi00} makes
the solution fully explicit.
% [AUTHOR: the $\gamma_1\to0^+$ reduction below is stated without proof. As
% $\gamma_1\to0^+$ both $\alpha=\mu/(\gamma_1 y)$ and $c=\lambda_1 y/\gamma_1\to\infty$;
% the claimed collapse of the Kummer ratio to the Model~$A$ rational form relies on a
% confluent-hypergeometric large-parameter asymptotic that should be derived or cited
% (e.g. DLMF 13.8). Flagging rather than silently asserting.]
In the limit $\gamma_1\to0^+$, the Kummer ratio
collapses to the rational expression $x^*(y)(1-y)/(x^*(y)-y)$ of Model~$A$, recovering
formula~\eqref{eq:A:fundamental}.
\end{proof}

% ─────────────────────────────────────────────────────────────────
\subsection{Mean queue lengths and mean waiting times}

\subsubsection*{Total queue length from the diagonal}

Setting $x=y=z$ in the reduced fundamental equation~\eqref{eq:B2:fundamental}, the
jockeying convection term $\gamma_1 xy(x-y)\partial_xP$ vanishes at $x=y=z$
(factor $(x-y)=0$), as does the boundary term $\mu(x-y)P_y(y)$.  The remaining
balance is identical to Model~$A$'s, so the same diagonal cancellation gives:
\begin{align}
    \begin{split}
        P(z,z) \;=\; \frac{\rho(1-\rho)}{1-\rho z},
    \end{split}
    \label{eq:B2:Pzz}
\end{align}
the same formula as~\eqref{eq:A:Pzz}.  Since $P(z,z)=\mathbb{E}[z^{N_1+N_2}]=\mathbb{E}[z^{N}]$
is the PGF of the total in-queue count $N$, the generating-function identity
$\tfrac{d}{dz}P(z,z)\big|_{z=1}=\mathbb{E}[N_1]+\mathbb{E}[N_2]$ holds;
differentiating by the quotient rule and evaluating at $z=1$:
\begin{align}
    \begin{split}
        \mathbb{E}[N_1]+\mathbb{E}[N_2]
        \;=\; \left.\frac{d}{dz}\,\frac{\rho(1-\rho)}{1-\rho z}\right|_{z=1}
        \;=\; \frac{\rho(1-\rho)\cdot\rho}{(1-\rho z)^2}\bigg|_{z=1}
        \;=\; \frac{\rho^2(1-\rho)}{(1-\rho)^2}
        \;=\; \frac{\rho^2}{1-\rho}.
    \end{split}
    \label{eq:B2:EN}
\end{align}
One-way jockeying does not alter the total mean queue length relative to
Model~$A$, since the diagonal PGF is unchanged.

\subsubsection*{Class-1 mean queue length}

We extract $\mathbb{E}[N_1]=\tfrac{\partial P}{\partial x}\big|_{(1,1)}$ by first
reducing the fundamental equation~\eqref{eq:B2:fundamental} at $y=1$, and then
differentiating the resulting marginal PGF at $x=1$.

\medskip\noindent\textit{Step~1: marginal PGF $P(x,1)$.}\enspace
Evaluating~\eqref{eq:B2:fundamental} at $y=1$, the bracket multiplying $P$ becomes
$(\lambda_1+\mu)x-\mu-\lambda_1 x^2=-\lambda_1(x-1)(x-1/\rho_1)$, the convection
coefficient is $\gamma_1 x(x-1)$, and the right-hand side reduces to $\mu(x-1)P_y(1)$
(the $\pi(0,0)$ term carries the factor $(1-y)=0$). Every term therefore shares the common
factor $(x-1)$; cancelling it leaves the first-order linear ODE in $x$:
\begin{align}
    \begin{split}
        -\lambda_1(x-1/\rho_1)\,P(x,1)
        \;+\; \gamma_1 x\,\frac{d}{dx}P(x,1)
        \;=\; \mu\,P_y(1).
    \end{split}
    \label{eq:B2:ode_y1}
\end{align}
The integrating factor for this ODE is $e^{-\lambda_1 x/\gamma_1}\,x^{\,\mu/\gamma_1}$
(noting $\lambda_1/\rho_1=\mu$), and integrating from $0$ to $x$ (the boundary term
vanishes since $0^{\mu/\gamma_1}=0$):
\begin{align}
    \begin{split}
        P(x,1)\,e^{-\lambda_1 x/\gamma_1}\,x^{\,\mu/\gamma_1}
        \;=\; \frac{\mu\,P_y(1)}{\gamma_1}
        \int_0^x e^{-\lambda_1 t/\gamma_1}\,t^{\,\mu/\gamma_1-1}\,dt.
    \end{split}
\end{align}
Rearranging and writing $\alpha=\mu/\gamma_1$ and
$\mathcal{J}(x)=\int_0^x e^{-\lambda_1 t/\gamma_1}\,t^{\,\alpha-1}\,dt$:
\begin{align}
    \begin{split}
        P(x,1)
        \;=\; \frac{\mu\,P_y(1)}{\gamma_1}
        \cdot e^{\lambda_1 x/\gamma_1}\,x^{-\alpha}\,\mathcal{J}(x).
    \end{split}
    \label{eq:B2:Px1}
\end{align}
The constant $P_y(1)$ follows from $P(1,1)=1-\pi_0=\rho$, which gives
$P_y(1)=\rho\gamma_1/\bigl(\mu\,e^{\lambda_1/\gamma_1}\,\mathcal{J}(1)\bigr)$.

\medskip\noindent\textit{Step~2: differentiate $P(x,1)$ at $x=1$.}\enspace
Writing $P(x,1)=C\cdot f(x)\cdot\mathcal{J}(x)$ with $C=\mu P_y(1)/\gamma_1$ and
$f(x)=e^{\lambda_1 x/\gamma_1}\,x^{-\alpha}$, the product rule together with
Leibniz's rule ($d\mathcal{J}/dx = e^{-\lambda_1 x/\gamma_1}\,x^{\alpha-1}$, the integrand
at the upper limit) gives:
\begin{align}
    \begin{split}
        \frac{d}{dx}P(x,1)
        &\;=\; C\Bigl[f'(x)\,\mathcal{J}(x)+f(x)\cdot e^{-\lambda_1 x/\gamma_1}\,x^{\alpha-1}\Bigr].
    \end{split}
\end{align}
The logarithmic derivative $f'(x)/f(x)=\lambda_1/\gamma_1-\alpha/x=(\lambda_1 x-\mu)/(\gamma_1 x)$,
and the second product term simplifies as
$f(x)\,e^{-\lambda_1 x/\gamma_1}x^{\alpha-1}=x^{-1}$, so:
\begin{align}
    \begin{split}
        \frac{d}{dx}P(x,1)
        &\;=\; \frac{\mu\,P_y(1)}{\gamma_1}\left[
            e^{\lambda_1 x/\gamma_1}\,x^{-\alpha}\cdot\frac{\lambda_1 x-\mu}{\gamma_1 x}
            \cdot\mathcal{J}(x)
            \;+\; x^{-1}
        \right].
    \end{split}
\end{align}

\medskip\noindent\textit{Step~3: evaluate at $x=1$.}\enspace
\begin{align}
    \begin{split}
        \mathbb{E}[N_1]
        &\;=\; \left.\frac{d}{dx}P(x,1)\right|_{x=1} \\
        &\;=\; \frac{\mu\,P_y(1)}{\gamma_1}\left[
            e^{\lambda_1/\gamma_1}\cdot\frac{\lambda_1-\mu}{\gamma_1}
            \cdot\mathcal{J}(1)
            \;+\; 1
        \right],
    \end{split}
    \label{eq:B2:EN1}
\end{align}
where $\mathcal{J}(1)=\int_0^1 e^{-\lambda_1 t/\gamma_1}\,t^{\,\mu/\gamma_1-1}\,dt
=(\gamma_1/\lambda_1)^{\mu/\gamma_1}\,\gamma(\mu/\gamma_1,\,\lambda_1/\gamma_1)$
and $\gamma(\cdot,\cdot)$ is the lower incomplete gamma function.  The
factor $(\lambda_1-\mu)<0$ (since $\rho_1=\lambda_1/\mu<1$), confirming that
$e^{\lambda_1/\gamma_1}(\lambda_1-\mu)\mathcal{J}(1)/\gamma_1 + 1 > 0$.
No elementary simplification of~\eqref{eq:B2:EN1} is known; the expression
is evaluated numerically.  As $\gamma_1\to0^+$ the
Model~$A$ value~\eqref{eq:A:EN1} is recovered.

\subsubsection*{Class-2 mean queue length by subtraction}

\begin{align}
    \begin{split}
        \mathbb{E}[N_2]
        \;=\; \bigl(\mathbb{E}[N_1]+\mathbb{E}[N_2]\bigr)-\mathbb{E}[N_1]
        \;=\; \frac{\rho^2}{1-\rho}-\mathbb{E}[N_1].
    \end{split}
    \label{eq:B2:EN2}
\end{align}
Since jockeying defections reduce the class-1 queue,
$\mathbb{E}[N_1]$ is smaller than the Model~$A$ value $\rho\rho_1/(1-\rho_1)$,
so $\mathbb{E}[N_2]$ correspondingly exceeds the Model~$A$ value.

\subsubsection*{Mean waiting times via Little's Law}

For Class~1, the entry rate into Queue~1 is exactly $\lambda_1$: every Class-1 arrival
joins Queue~1, whether it later jockeys to Queue~2 or enters service. \textbf{Little's
Law} applied to the Queue-1 waiting line therefore gives the genuine mean time in Queue~1,
\begin{align}
    \begin{split}
        \mathbb{E}[W_1] \;=\; \frac{\mathbb{E}[N_1]}{\lambda_1}.
    \end{split}
    \label{eq:B2:EW}
\end{align}

% [AUTHOR: throughput for E[W_2]. Unlike Models A and C_2, Queue 2 in Model B_2 has TWO
% input streams: the Poisson arrivals at rate lambda_2 AND the jockeying transfers from
% Queue 1, whose long-run rate is gamma_1 * E[N_1]. Little's Law applied to the Queue-2
% buffer reads E[N_2] = (lambda_2 + gamma_1*E[N_1]) * E[W_2], so the mean time in Queue 2
% is E[N_2]/(lambda_2 + gamma_1*E[N_1]), NOT E[N_2]/lambda_2. The naive E[N_2]/lambda_2 is
% correct only if E[W_2] is DEFINED as the Little ratio per Class-2 arrival rather than as
% a sojourn time. Please confirm the intended definition of W_2 and fix the denominator
% (or the prose) accordingly. The relation below uses the sojourn-time convention.]
For Class~2, Queue~2 receives not only the Poisson arrivals at rate $\lambda_2$ but also
the jockeying transfers from Queue~1, whose long-run rate is $\gamma_1\,\mathbb{E}[N_1]$.
\textbf{Little's Law} applied to the Queue-2 buffer therefore reads
$\mathbb{E}[N_2]=\bigl(\lambda_2+\gamma_1\,\mathbb{E}[N_1]\bigr)\,\mathbb{E}[W_2]$, so the
mean sojourn in Queue~2 is
\begin{align}
    \begin{split}
        \mathbb{E}[W_2] \;=\; \frac{\mathbb{E}[N_2]}{\lambda_2+\gamma_1\,\mathbb{E}[N_1]},
    \end{split}
    \label{eq:B2:EW2}
\end{align}
which reduces to $\mathbb{E}[N_2]/\lambda_2$ only in the no-jockeying limit
$\gamma_1\to0^+$.  As $\gamma_1\to0^+$ both~\eqref{eq:B2:EW} and~\eqref{eq:B2:EW2}
recover the Model~$A$ values~\eqref{eq:A:EW1}--\eqref{eq:A:EW2}.
